\section{Continuous maps and sheaves of modules}
\label{section-presheaves-modules-functorial}

\noindent
The case of sheaves of modules is more complicated.
The reason is that the natural setting for defining
the pullback and pushforward functors, is the setting
of ringed spaces, which we will define below. First
we state a few obvious lemmas.

\begin{lemma}
\label{lemma-pushforward-presheaf-module}
Let $f : X \to Y$ be a continuous map of topological spaces.
Let $\mathcal{O}$ be a presheaf of rings on $X$. Let
$\mathcal{F}$ be a presheaf of $\mathcal{O}$-modules.
There is a natural map of underlying presheaves of sets
$$
f_*\mathcal{O} \times f_*\mathcal{F}
\longrightarrow
f_*\mathcal{F}
$$
which turns $f_*\mathcal{F}$ into a presheaf of
$f_*\mathcal{O}$-modules. This construction is
functorial in $\mathcal{F}$.
\end{lemma}

\begin{proof}
Let $V \subset Y$ is open. We define the map of the lemma
to be the map
$$
f_*\mathcal{O}(V) \times f_*\mathcal{F}(V)
=
\mathcal{O}(f^{-1}V) \times \mathcal{F}(f^{-1}V)
\to
\mathcal{F}(f^{-1}V)
=
f_*\mathcal{F}(V).
$$
Here the arrow in the middle is the multiplication map on $X$.
We leave it to the reader to see this is compatible with
restriction mappings and defines a structure of
$f_*\mathcal{O}$-module on $f_*\mathcal{F}$.
\end{proof}

\begin{lemma}
\label{lemma-pullback-presheaf-module}
Let $f : X \to Y$ be a continuous map of topological spaces.
Let $\mathcal{O}$ be a presheaf of rings on $Y$. Let
$\mathcal{G}$ be a presheaf of $\mathcal{O}$-modules.
There is a natural map of underlying presheaves of sets
$$
f_p\mathcal{O} \times f_p\mathcal{G}
\longrightarrow
f_p\mathcal{G}
$$
which turns $f_p\mathcal{G}$ into a presheaf of $f_p\mathcal{O}$-modules.
This construction is functorial in $\mathcal{G}$.
\end{lemma}

\begin{proof}
Let $U \subset X$ is open. We define the map of the lemma
to be the map
\begin{eqnarray*}
f_p\mathcal{O}(U) \times f_p\mathcal{G}(U)
& = &
\colim_{f(U) \subset V} \mathcal{O}(V)
\times
\colim_{f(U) \subset V} \mathcal{G}(V) \\
& = &
\colim_{f(U) \subset V} (\mathcal{O}(V)\times \mathcal{G}(V)) \\
& \to &
\colim_{f(U) \subset V} \mathcal{G}(V) \\
& = &
f_p\mathcal{G}(U).
\end{eqnarray*}
Here the arrow in the middle is the multiplication map on $Y$.
The second equality holds because directed colimits commute
with finite limits, see
Categories, Lemma \ref{categories-lemma-directed-commutes}.
We leave it to the reader to see this is compatible with
restriction mappings and defines a structure of
$f_p\mathcal{O}$-module on $f_p\mathcal{G}$.
\end{proof}

\noindent
Let $f : X \to Y$ be a continuous map.
Let $\mathcal{O}_X$ be a presheaf of rings on $X$ and
let $\mathcal{O}_Y$ be a presheaf of rings on $Y$.
So at the moment we have defined functors
\begin{eqnarray*}
f_* : \textit{PMod}(\mathcal{O}_X) &
\longrightarrow &
\textit{PMod}(f_*\mathcal{O}_X) \\
f_p : \textit{PMod}(\mathcal{O}_Y) &
\longrightarrow &
\textit{PMod}(f_p\mathcal{O}_Y)
\end{eqnarray*}
These satisfy some compatibilities as follows.

\begin{lemma}
\label{lemma-adjoint-push-pull-presheaves-modules}
Let $f : X \to Y$ be a continuous map of topological spaces.
Let $\mathcal{O}$ be a presheaf of rings on $Y$.
Let $\mathcal{G}$ be a presheaf of $\mathcal{O}$-modules.
Let $\mathcal{F}$ be a presheaf of $f_p\mathcal{O}$-modules.
Then
$$
\Mor_{\textit{PMod}(f_p\mathcal{O})}(f_p\mathcal{G}, \mathcal{F})
=
\Mor_{\textit{PMod}(\mathcal{O})}(\mathcal{G}, f_*\mathcal{F}).
$$
Here we use
Lemmas \ref{lemma-pullback-presheaf-module}
and \ref{lemma-pushforward-presheaf-module}, and we think of
$f_*\mathcal{F}$ as an $\mathcal{O}$-module via the map
$i_\mathcal{O} : \mathcal{O} \to f_*f_p\mathcal{O}$
(defined first in the proof of Lemma \ref{lemma-pullback-presheaves}).
\end{lemma}

\begin{proof}
Note that we have
$$
\Mor_{\textit{PAb}(X)}(f_p\mathcal{G}, \mathcal{F})
=
\Mor_{\textit{PAb}(Y)}(\mathcal{G}, f_*\mathcal{F}).
$$
according to Section \ref{section-abelian-presheaves-functorial}.
So what we have to prove is that under this correspondence, the
subsets of module maps correspond. In addition, the correspondence
is determined by the rule
$$
(\psi : f_p\mathcal{G} \to \mathcal{F})
\longmapsto
(f_*\psi \circ i_\mathcal{G} :
\mathcal{G} \to f_* \mathcal{F})
$$
and in the other direction by the rule
$$
(\varphi : \mathcal{G} \to f_* \mathcal{F})
\longmapsto
(c_\mathcal{F} \circ f_p\varphi : f_p\mathcal{G} \to \mathcal{F})
$$
where $i_\mathcal{G}$ and $c_\mathcal{F}$ are as in
Section \ref{section-abelian-presheaves-functorial}.
Hence, using the functoriality of $f_*$ and $f_p$ we see that
it suffices to check that the maps
$i_\mathcal{G} : \mathcal{G} \to f_* f_p \mathcal{G}$ and
$c_\mathcal{F} : f_p f_* \mathcal{F} \to \mathcal{F}$
are compatible with module structures, which we leave to the reader.
\end{proof}

\begin{lemma}
\label{lemma-adjoint-pull-push-presheaves-modules}
Let $f : X \to Y$ be a continuous map of topological spaces.
Let $\mathcal{O}$ be a presheaf of rings on $X$.
Let $\mathcal{F}$ be a presheaf of $\mathcal{O}$-modules.
Let $\mathcal{G}$ be a presheaf of $f_*\mathcal{O}$-modules.
Then
$$
\Mor_{\textit{PMod}(\mathcal{O})}(
\mathcal{O} \otimes_{p, f_pf_*\mathcal{O}} f_p\mathcal{G}, \mathcal{F})
=
\Mor_{\textit{PMod}(f_*\mathcal{O})}(\mathcal{G}, f_*\mathcal{F}).
$$
Here we use
Lemmas \ref{lemma-pullback-presheaf-module}
and \ref{lemma-pushforward-presheaf-module}, and we use
the map $c_\mathcal{O} : f_pf_*\mathcal{O} \to \mathcal{O}$
in the definition of the tensor product.
\end{lemma}

\begin{proof}
This follows from the equalities
\begin{eqnarray*}
\Mor_{\textit{PMod}(\mathcal{O})}(
\mathcal{O} \otimes_{p, f_pf_*\mathcal{O}} f_p\mathcal{G}, \mathcal{F})
& = &
\Mor_{\textit{PMod}(f_pf_*\mathcal{O})}(
f_p\mathcal{G}, \mathcal{F}_{f_pf_*\mathcal{O}}) \\
& = &
\Mor_{\textit{PMod}(f_*\mathcal{O})}(\mathcal{G},
f_*(\mathcal{F}_{f_pf_*\mathcal{O}})) \\
& = &
\Mor_{\textit{PMod}(f_*\mathcal{O})}(\mathcal{G}, f_*\mathcal{F}).
\end{eqnarray*}
The first equality is
Lemma \ref{lemma-adjointness-tensor-restrict-presheaves}.
The second equality is
Lemma \ref{lemma-adjoint-push-pull-presheaves-modules}.
The third equality is given by the equality
$f_*(\mathcal{F}_{f_pf_*\mathcal{O}}) = f_*\mathcal{F}$
of abelian sheaves which is $f_*\mathcal{O}$-linear. Namely,
$\text{id}_{f_*\mathcal{O}}$ corresponds to $c_\mathcal{O}$
under the adjunction described in the proof of
Lemma \ref{lemma-pullback-presheaves}
and thus
$\text{id}_{f_*\mathcal{O}} = f_*c_\mathcal{O} \circ i_{f_*\mathcal{O}}$.
\end{proof}

\begin{lemma}
\label{lemma-pushforward-module}
Let $f : X \to Y$ be a continuous map of topological spaces.
Let $\mathcal{O}$ be a sheaf of rings on $X$. Let
$\mathcal{F}$ be a sheaf of $\mathcal{O}$-modules.
The pushforward $f_*\mathcal{F}$, as defined in
Lemma \ref{lemma-pushforward-presheaf-module}
is a sheaf of $f_*\mathcal{O}$-modules.
\end{lemma}

\begin{proof}
Obvious from the definition and Lemma \ref{lemma-pushforward-sheaf}.
\end{proof}

\begin{lemma}
\label{lemma-pullback-module}
Let $f : X \to Y$ be a continuous map of topological spaces.
Let $\mathcal{O}$ be a sheaf of rings on $Y$. Let
$\mathcal{G}$ be a sheaf of $\mathcal{O}$-modules.
There is a natural map of underlying presheaves of sets
$$
f^{-1}\mathcal{O} \times f^{-1}\mathcal{G}
\longrightarrow
f^{-1}\mathcal{G}
$$
which turns $f^{-1}\mathcal{G}$ into a
sheaf of $f^{-1}\mathcal{O}$-modules.
\end{lemma}

\begin{proof}
Recall that $f^{-1}$ is defined as the composition of the
functor $f_p$ and sheafification. Thus the lemma
is a combination of Lemma \ref{lemma-pullback-presheaf-module}
and Lemma \ref{lemma-sheafification-presheaf-modules}.
\end{proof}

\noindent
Let $f : X \to Y$ be a continuous map.
Let $\mathcal{O}_X$ be a sheaf of rings on $X$ and
let $\mathcal{O}_Y$ be a sheaf of rings on $Y$.
So now we have defined functors
\begin{eqnarray*}
f_* : \textit{Mod}(\mathcal{O}_X) &
\longrightarrow &
\textit{Mod}(f_*\mathcal{O}_X) \\
f^{-1} : \textit{Mod}(\mathcal{O}_Y) &
\longrightarrow &
\textit{Mod}(f^{-1}\mathcal{O}_Y)
\end{eqnarray*}
These satisfy some compatibilities as follows.

\begin{lemma}
\label{lemma-adjoint-push-pull-modules}
Let $f : X \to Y$ be a continuous map of topological spaces.
Let $\mathcal{O}$ be a sheaf of rings on $Y$.
Let $\mathcal{G}$ be a sheaf of $\mathcal{O}$-modules.
Let $\mathcal{F}$ be a sheaf of $f^{-1}\mathcal{O}$-modules.
Then
$$
\Mor_{\textit{Mod}(f^{-1}\mathcal{O})}(f^{-1}\mathcal{G}, \mathcal{F})
=
\Mor_{\textit{Mod}(\mathcal{O})}(\mathcal{G}, f_*\mathcal{F}).
$$
Here we use
Lemmas \ref{lemma-pullback-module}
and \ref{lemma-pushforward-module}, and we think of
$f_*\mathcal{F}$ as an $\mathcal{O}$-module by restriction via
$\mathcal{O} \to f_*f^{-1}\mathcal{O}$.
\end{lemma}

\begin{proof}
Argue by the equalities
\begin{eqnarray*}
\Mor_{\textit{Mod}(f^{-1}\mathcal{O})}(f^{-1}\mathcal{G}, \mathcal{F})
& = &
\Mor_{\textit{Mod}(f_p\mathcal{O})}(f_p\mathcal{G}, \mathcal{F}) \\
& = &
\Mor_{\textit{Mod}(\mathcal{O})}(\mathcal{G}, f_*\mathcal{F}).
\end{eqnarray*}
where the second is
Lemmas \ref{lemma-adjoint-push-pull-presheaves-modules}
and the first is by Lemma \ref{lemma-sheafification-presheaf-modules}.
\end{proof}

\begin{lemma}
\label{lemma-adjoint-pull-push-modules}
Let $f : X \to Y$ be a continuous map of topological spaces.
Let $\mathcal{O}$ be a sheaf of rings on $X$.
Let $\mathcal{F}$ be a sheaf of $\mathcal{O}$-modules.
Let $\mathcal{G}$ be a sheaf of $f_*\mathcal{O}$-modules.
Then
$$
\Mor_{\textit{Mod}(\mathcal{O})}(
\mathcal{O} \otimes_{f^{-1}f_*\mathcal{O}} f^{-1}\mathcal{G}, \mathcal{F})
=
\Mor_{\textit{Mod}(f_*\mathcal{O})}(\mathcal{G}, f_*\mathcal{F}).
$$
Here we use
Lemmas \ref{lemma-pullback-module}
and \ref{lemma-pushforward-module}, and we use
the canonical map $f^{-1}f_*\mathcal{O} \to \mathcal{O}$
in the definition of the tensor product.
\end{lemma}

\begin{proof}
This follows from the equalities
\begin{eqnarray*}
\Mor_{\textit{Mod}(\mathcal{O})}(
\mathcal{O} \otimes_{f^{-1}f_*\mathcal{O}} f^{-1}\mathcal{G}, \mathcal{F})
& = &
\Mor_{\textit{Mod}(f^{-1}f_*\mathcal{O})}(
f^{-1}\mathcal{G}, \mathcal{F}_{f^{-1}f_*\mathcal{O}}) \\
& = &
\Mor_{\textit{Mod}(f_*\mathcal{O})}(\mathcal{G}, f_*\mathcal{F}).
\end{eqnarray*}
which are a combination of
Lemma \ref{lemma-adjointness-tensor-restrict}
and \ref{lemma-adjoint-push-pull-modules}.
\end{proof}

\noindent
Let $f : X \to Y$ be a continuous map.
Let $\mathcal{O}_X$ be a (pre)sheaf of rings on $X$ and
let $\mathcal{O}_Y$ be a (pre)sheaf of rings on $Y$.
So at the moment we have defined functors
\begin{eqnarray*}
f_* : \textit{PMod}(\mathcal{O}_X) &
\longrightarrow &
\textit{PMod}(f_*\mathcal{O}_X) \\
f_* : \textit{Mod}(\mathcal{O}_X) &
\longrightarrow &
\textit{Mod}(f_*\mathcal{O}_X) \\
f_p : \textit{PMod}(\mathcal{O}_Y) &
\longrightarrow &
\textit{PMod}(f_p\mathcal{O}_Y) \\
f^{-1} : \textit{Mod}(\mathcal{O}_Y) &
\longrightarrow &
\textit{Mod}(f^{-1}\mathcal{O}_Y)
\end{eqnarray*}
Clearly, usually the pair of functors $(f_*, f^{-1})$
on sheaves of modules are not adjoint, because their target categories
do not match. Namely, as we saw above, it works only if by some miracle the
sheaves of rings $\mathcal{O}_X, \mathcal{O}_Y$ satisfy the
relations $\mathcal{O}_X = f^{-1}\mathcal{O}_Y$ and
$\mathcal{O}_Y = f_*\mathcal{O}_X$. This is almost never
true in practice. We interrupt the discussion to define
the correct notion of morphism for which a suitable adjoint
pair of functors on sheaves of modules exists.
